% GENERATED FILE. Edit canonical lessons, then run npm run generate:books.
% canonical-source-hash: fnv1a64:e1c8dcf93516e680
% canonical-lessons: KA-C246-uluku, KA-C246-ajirna, KA-C246-hunnu, KA-C246-byandej, KA-C246-madhumeha

\chapter{Sprain, Indigestion, Sore, Bandage, Diabetes}
\label{ch:ka-sprain-indigestion-sore-bandage-diabetes}
\hlchaptermodality{246}

\begin{chapteropening}
\textbf{By the end of this chapter:} \emph{I can say a sprain, indigestion, a sore, a bandage and diabetes.}

Complete the last lesson of chapter 246: I can say a sprain, indigestion, a sore, a bandage and diabetes.
\end{chapteropening}

\section[\texorpdfstring{\kn{ಉಳುಕು} (uḷuku)}{ಉಳುಕು (uḷuku)}]{\kn{ಉಳುಕು} (uḷuku) — a sprain}

\label{lesson:KA-C246-uluku}



\begin{quote}
Before the new one: say the Kannada for a pimple, then the Kannada for medical treatment.
\end{quote}

\subsection*{You'll want to know: \kn{ಉಳುಕು}}
\textbf{\kn{ಉಳುಕು}} — \emph{uḷuku} — \textquotedblleft{}a sprain\textquotedblright{}.

\begin{grammarlens}[title={What you've built}]
One more word for this chapter.
\end{grammarlens}

\subsection*{Your turn}
\begin{itemize}
\raggedright
  \item \emph{Say it:} \emph{uḷuku}
  \item \emph{Say it:} \emph{uḷuku}, once more
  \item \emph{Say it:} say \emph{uḷuku}
  \item \emph{Recall:} say the Kannada for a pimple, then the Kannada for medical treatment, then say \emph{uḷuku} again
\end{itemize}

\begin{tcolorbox}[breakable,colback=teal!4,colframe=teal!35!black,title={Before you move on}]
What does \kn{ಉಳುಕು} mean? (\textquotedblleft{}A sprain\textquotedblright{}.) Say it once more.
\end{tcolorbox}
\section[\texorpdfstring{\kn{ಅಜೀರ್ಣ} (ajīrṇa)}{ಅಜೀರ್ಣ (ajīrṇa)}]{\kn{ಅಜೀರ್ಣ} (ajīrṇa) — indigestion}

\label{lesson:KA-C246-ajirna}



\begin{quote}
Before the new one: say the Kannada for medical treatment, then the Kannada for a sprain.
\end{quote}

\subsection*{You'll want to know: \kn{ಅಜೀರ್ಣ}}
\textbf{\kn{ಅಜೀರ್ಣ}} — \emph{ajīrṇa} — \textquotedblleft{}indigestion\textquotedblright{}.

\begin{grammarlens}[title={What you've built}]
One more word for this chapter.
\end{grammarlens}

\subsection*{Your turn}
\begin{itemize}
\raggedright
  \item \emph{Say it:} \emph{ajīrṇa}
  \item \emph{Say it:} \emph{ajīrṇa}, once more
  \item \emph{Say it:} say \emph{ajīrṇa}
  \item \emph{Recall:} say the Kannada for medical treatment, then the Kannada for a sprain, then say \emph{ajīrṇa} again
\end{itemize}

\begin{tcolorbox}[breakable,colback=teal!4,colframe=teal!35!black,title={Before you move on}]
What does \kn{ಅಜೀರ್ಣ} mean? (\textquotedblleft{}Indigestion\textquotedblright{}.) Say it once more.
\end{tcolorbox}
\section[\texorpdfstring{\kn{ಹುಣ್ಣು} (huṇṇu)}{ಹುಣ್ಣು (huṇṇu)}]{\kn{ಹುಣ್ಣು} (huṇṇu) — a sore, an ulcer}

\label{lesson:KA-C246-hunnu}



\begin{quote}
Before the new one: say the Kannada for a sprain, then the Kannada for indigestion.
\end{quote}

\subsection*{You'll want to know: \kn{ಹುಣ್ಣು}}
\textbf{\kn{ಹುಣ್ಣು}} — \emph{huṇṇu} — \textquotedblleft{}a sore, an ulcer\textquotedblright{}.

\begin{grammarlens}[title={What you've built}]
One more word for this chapter.
\end{grammarlens}

\subsection*{Your turn}
\begin{itemize}
\raggedright
  \item \emph{Say it:} \emph{huṇṇu}
  \item \emph{Say it:} \emph{huṇṇu}, once more
  \item \emph{Say it:} say \emph{huṇṇu}
  \item \emph{Recall:} say the Kannada for a sprain, then the Kannada for indigestion, then say \emph{huṇṇu} again
\end{itemize}

\begin{tcolorbox}[breakable,colback=teal!4,colframe=teal!35!black,title={Before you move on}]
What does \kn{ಹುಣ್ಣು} mean? (\textquotedblleft{}A sore, an ulcer\textquotedblright{}.) Say it once more.
\end{tcolorbox}
\section[\texorpdfstring{\kn{ಬ್ಯಾಂಡೇಜ್} (byāṇḍēj)}{ಬ್ಯಾಂಡೇಜ್ (byāṇḍēj)}]{\kn{ಬ್ಯಾಂಡೇಜ್} (byāṇḍēj) — a bandage}

\label{lesson:KA-C246-byandej}



\begin{quote}
Before the new one: say the Kannada for indigestion, then the Kannada for a sore.
\end{quote}

\subsection*{You'll want to know: \kn{ಬ್ಯಾಂಡೇಜ್}}
\textbf{\kn{ಬ್ಯಾಂಡೇಜ್}} — \emph{byāṇḍēj} — \textquotedblleft{}a bandage\textquotedblright{}.

\begin{grammarlens}[title={What you've built}]
One more word for this chapter.
\end{grammarlens}

\subsection*{Your turn}
\begin{itemize}
\raggedright
  \item \emph{Say it:} \emph{byāṇḍēj}
  \item \emph{Say it:} \emph{byāṇḍēj}, once more
  \item \emph{Say it:} say \emph{byāṇḍēj}
  \item \emph{Recall:} say the Kannada for indigestion, then the Kannada for a sore, then say \emph{byāṇḍēj} again
\end{itemize}

\begin{tcolorbox}[breakable,colback=teal!4,colframe=teal!35!black,title={Before you move on}]
What does \kn{ಬ್ಯಾಂಡೇಜ್} mean? (\textquotedblleft{}A bandage\textquotedblright{}.) Say it once more.
\end{tcolorbox}
\section[\texorpdfstring{\kn{ಮಧುಮೇಹ} (madhumēha)}{ಮಧುಮೇಹ (madhumēha)}]{\kn{ಮಧುಮೇಹ} (madhumēha) — diabetes}

\label{lesson:KA-C246-madhumeha}



\begin{quote}
Before the new one: say the Kannada for a sore, then the Kannada for a bandage.
\end{quote}

\subsection*{You'll want to know: \kn{ಮಧುಮೇಹ}}
\textbf{\kn{ಮಧುಮೇಹ}} — \emph{madhumēha} — \textquotedblleft{}diabetes\textquotedblright{}.

\begin{grammarlens}[title={What you've built}]
One more word for this chapter.
\end{grammarlens}

\subsection*{Your turn}
\begin{itemize}
\raggedright
  \item \emph{Say it:} \emph{madhumēha}
  \item \emph{Say it:} \emph{madhumēha}, once more
  \item \emph{Say it:} say \emph{madhumēha}
  \item \emph{Recall:} say the Kannada for a sore, then the Kannada for a bandage, then say \emph{madhumēha} again
\end{itemize}

\begin{tcolorbox}[breakable,colback=teal!4,colframe=teal!35!black,title={Before you move on}]
What does \kn{ಮಧುಮೇಹ} mean? (\textquotedblleft{}Diabetes\textquotedblright{}.) Say it once more.
\end{tcolorbox}
